\section{Discussion}
\label{sec:discussion}

\subsection{Architectural commitments}

Three architectural commitments shape the construction.

\paragraph{Dependency, not syntax, is the carrier of meaning.}
The semantic owner is the event graph, not the source program.  A
consequence is that two source programs that elaborate to the same
event graph are strategically indistinguishable -- by construction,
not by an after-the-fact equivalence proof.  This frees the surface
language to add ergonomic constructors (simultaneous yields,
nullable yields, future handlers) without changing the strategic
semantics; lowering is the only translation that can change the
event-graph shape.

\paragraph{Configurations are extensional.}
A configuration records which nodes have produced which results;
schedules are not part of the state.  This is what makes the
frontier diamond (\Cref{thm:diamond}) definitional rather than a
quotient construction.  The cost is that an interpreter that wants
to walk a specific schedule must supply that schedule explicitly as an
event law, batch witness, checkpoint policy, or runtime trace law.  The
graph itself does not privilege a source linearization.

\paragraph{Visibility is bound, not used.}
Every binding has a visibility tag; expression typing then
projects the context down to public or hidden as required.  The
discipline is owner-or-all: there is no group-visibility
constructor, and reveal-like operations introduce a fresh public
alias rather than mutating visibility.  Both decisions keep the
view algebra small (two erasures and a per-player view) and let
guard-uses-only-style predicates be about \emph{dependency
sets}, not about typing rules.

\subsection{Limitations}
\label{sec:disc:limits}

\paragraph{Finiteness.}  Player sets, value domains, programs, event
graphs, and reachable strategy carriers are finite.  Distributions
are finite-support PMFs.  Continuous types and measure-theoretic
kernels~\citep{Giry1982,Kozen1981} are out of scope; the bridge
from $\FWeight{\cdot}$ to mathlib's $\PMF{\cdot}$ is one-way and
takes a normalization proof, with no measure-theoretic continuation.

\paragraph{No quit handlers.}  The surface language deliberately
omits user-level quit handlers.  A strategic quit changes the
action space and the control-flow shape, neither of which is a
value-level effect; supporting it would require adding a
prerequisite-cancellation rule to the event graph and revisiting
the frontier-read invariants.  Nullable yields are the present
substitute for ``opt out of revealing.''

\paragraph{Solution-concept inventory.}  The strategic API exposes
Nash, dominance, correlated and coarse correlated equilibria,
potential-game predicates, zero-sum/team predicates, social welfare,
and Pareto efficiency.  Refinement-based equilibrium concepts
(subgame-perfect equilibria, sequential equilibria,
trembling-hand~\citep{Selten1975}) are not yet exposed because the
event-graph view does not maintain a sub-game decomposition; this
would require an additional presentation or subgame API rather than
only another predicate on the existing kernel games.

\paragraph{Presentation boundaries.}  The Vegas-facing pure and
PMF-behavioral strategy carriers are native round-view carriers over the
event-graph machine.  The bounded FOSG view is the standard
imperfect-information presentation, and the EFG layer is the serialized
tree presentation obtained from it.  The public Kuhn theorems are stated
and proved at the native completed-frontier level.  FOSG and EFG
adequacy results can be composed with those theorems when an exported
presentation is needed; they do not own the checked-program Kuhn proof.

\paragraph{Public samples only.}  The source $\Sample$ constructor
draws a public coin.  It is suitable for public randomness, but it
does not create private Harsanyi types or hidden chance moves.
Those require either a new hidden/private sample primitive or an
explicit external type parameter that a player commits to under the
visibility discipline.

\paragraph{Solver export.}  The finite pure kernel-game export object
contains player/strategy enumerations and kernel/utility oracles.  It is
deliberately narrower than a complete solver interchange format, and the
behavioral carrier is not exported as a finite strategy list.

\paragraph{Decidability.}  The static obligations of
\S\ref{sec:source:wf} -- contextual well-formedness, fresh
bindings, reveal completeness, and distribution normalization --
are decidable.  Legality is decidable when value domains are
finite, which is the standing assumption.  The Nash predicate and friends are quantifier-bounded
over finite carriers, so they are decidable in principle; whether
they are computable in useful time is a separate question
\citep{DaskalakisGoldbergPapadimitriou2009}.

\subsection{Natural extensions}
\label{sec:disc:ext}

\paragraph{Mechanism design under epistemic types.}
Bayesian games with type
spaces~\citep{Harsanyi1967,MyersonGameTheory} are close to the
intended scope but need one more primitive: private type draws.
The current $\Sample$ layer is public, so it cannot by itself
generate hidden per-player types.  A mechanism-design library could
be packaged as a checked \TT{WFProgram} builder over a type-space
schema once the source language has either hidden samples or an
external typed-input interface.

\paragraph{Concurrent and partial-information protocols at scale.}
The event graph is naturally an event structure in the sense of
\citet{Winskel1986EventStructures,NielsenPlotkinWinskel1981}.
Concurrent-game-semantic
techniques~\citep{RidemanWinskel2011Concurrent,
CastellanYoshidaSessionGame,
AbramskyJagadeesan1994,HylandOng2000PCF}
and session-typed multi-party
protocols~\citep{HondaYoshidaCarbone2008MPST} suggest two
extensions: typed channel structure beyond the
single-public-write SSA convention, and a richer notion of session
that tracks delegated authority.  Both are layerable on top of the
present event-graph carrier.

\paragraph{Cryptographic backends.}
The commit/reveal pattern is operationally a cryptographic
commitment scheme~\citep{Blum1983CoinFlipping,
GoldwasserMicaliRivest1988,Pedersen1991Commitment}.  The current
construction abstracts the commitment as a sealed-then-revealed
binding; a refinement layer that models hash-based or
homomorphic-commitment implementations would be a natural backend
under stochastic refinement, with the cryptographic soundness
property appearing as a side condition on $\refines_s$.

\paragraph{Compositional structure.}
Every checked program produces a kernel game, and morphisms between
kernel games are EU-preserving.  This suggests a categorical
structure on checked programs analogous to compositional game
theory~\citep{GhaniHedgesWinschel2018CGT,Hedges2016CGT};
making the composition explicit (parallel,
sequential, and conditional combinators on checked programs)
would let large protocols be assembled from small ones with
machine-checked strategic-structure preservation.

\paragraph{Concurrent-strategy realization.}
The frontier Kuhn bridge gives the realization half of the classical
mixed-vs-behavioral correspondence.  In a concurrent-strategy
setting (where a player can commit to a function of \emph{joint}
information across simultaneous events), the corresponding
realization theorem is more delicate, and the formalism naturally
intersects logical accounts of partial information
\citep{HintikkaSandu1989,MannSanduSevenster2011IFL}.

\subsection{Outlook}

The construction settles a small but central question: when the
strategic content of a finite multi-party protocol is mechanically
extracted from the program, what is the right primary semantic
object?  The answer it offers is the dependency graph of events,
with extensional configurations and a checked frontier.  The
strategic-form, FOSG, EFG, and backend views are then
\emph{presentations} of this object that come with their own
adequacy and transport theorems
(payoff-vector agreement, native frontier Kuhn realization, EFG presentation
adequacy, refinement transport).

The deeper benefit is that future extensions -- compositional
combinators, larger surface notation, richer cryptographic backends,
mechanism-design packaging -- can all be added as new
presentations or new lowerings without touching the central
event-graph theorems.  Adequacy at the boundary is a checked
theorem; the inside of the boundary is free.
